\section{Ejemplos de prueba}

\begin{longtable}{>{\raggedright\arraybackslash}p{0.32\linewidth}>{\raggedright\arraybackslash}p{0.22\linewidth}>{\raggedright\arraybackslash}p{0.36\linewidth}}
\toprule
\textbf{Comando} & \textbf{Estado} & \textbf{Observación} \\
\midrule
\texttt{make clean \&\& make} & Verificado local y Pampero & Compila
\texttt{bin/server} y \texttt{bin/client} sin warnings visibles. \\
\texttt{make test} & Verificado local y Pampero & Unitarios C del parser,
usuarios, métricas, logger, selector bloqueante y relay. \\
\texttt{make check PORT=12080} & Verificado local y Pampero & Flujos completos
de negociación, autenticación, conexión, relay, DNS, métricas y PMC. \\
\texttt{make valgrind PORT=...} & Verificado en Pampero & Valgrind con tráfico:
0 errores, 0 leaks definitivos/indirectos/posibles. \\
\texttt{strace} sobre relay HTTP & Verificado en Pampero & En ambos sentidos se
observó \texttt{recvfrom} $\rightarrow$ \texttt{sendto} sin un
\texttt{pselect6} intermedio. \\
\texttt{test/stress/run\_stress.py} & Corrida local acotada & Genera CSV y
figuras reproducibles para las combinaciones elegidas. \\
\bottomrule
\end{longtable}

\subsection{Corrida Pampero/Linux}

La corrida remota se ejecutó en \texttt{pampero.it.itba.edu.ar}, Linux
\texttt{7.0.9-arch1-1}, GCC \texttt{16.1.1 20260625}. El build
\texttt{make clean \&\& make} terminó con \texttt{BUILD OK}.

Los unitarios reportaron todos los casos OK y cubrieron:
\begin{itemize}[nosep]
  \item negociación, autenticación y parseo de pedidos fragmentados;
  \item conexión, relay bidireccional, cierres parciales y contrapresión;
  \item notificaciones del selector, métricas y registro de accesos.
\end{itemize}

Las integraciones terminaron sin fallas y ejercitaron los mismos comportamientos
sobre sockets reales. El único caso no ejecutado requería que el entorno
resolviera un dominio a varias direcciones para forzar el reintento; esa
precondición no estaba disponible en la corrida.

Valgrind con tráfico real reportó:
\begin{itemize}[nosep]
  \item \texttt{FILE DESCRIPTORS: 2 open (2 inherited) at exit};
  \item \texttt{definitely lost: 0 bytes}, \texttt{indirectly lost: 0 bytes} y
  \texttt{possibly lost: 0 bytes};
  \item \texttt{ERROR SUMMARY: 0 errors from 0 contexts}.
\end{itemize}
Los \texttt{5,658 bytes} \texttt{still reachable} se consideran alcanzables del
entorno/libc, no leaks definitivos del programa.

\subsection{Stress local}

El banco de stress se genera con \texttt{test/stress/run\_stress.py}. El script
levanta un origin controlado, ejecuta el servidor SOCKS5, consulta métricas por
PMC y escribe CSV bajo un directorio de salida. Los parámetros principales son:

\begin{itemize}[nosep]
  \item concurrencia: cantidad de conexiones simultáneas por caso;
  \item payload: bytes transferidos por conexión;
  \item modo de destino: \texttt{ipv4} o \texttt{fqdn};
  \item métricas recolectadas: conexiones OK/fallidas, bytes, wall time,
  throughput, tasa de conexiones, latencias p50/p95, REP no exitosos y deltas
  de métricas PMC.
\end{itemize}

La evidencia versionada corresponde a una corrida local acotada en macOS,
guardada en:
\begin{center}
\texttt{docs/stress/2026-07-05\_local-smoke\_v1/}
\end{center}

La corrida usó concurrencias 5 y 20, payloads 1024 y 8192 bytes, destinos IPv4 y
FQDN, y una repetición por combinación. También hubo una prueba local adicional de
500 conexiones con \texttt{conn\_ok=500} y \texttt{conn\_failed=0}. No se
presenta como máximo absoluto del servidor ni como barrido completo en Pampero.

\IfFileExists{../stress/2026-07-05_local-smoke_v1/2026-07-05_fig-throughput-vs-concurrency_v1.png}{
\begin{figure}[htbp]
  \centering
  \includegraphics[width=0.82\linewidth]{2026-07-05_fig-throughput-vs-concurrency_v1.png}
  \caption{Throughput agregado de la corrida local por concurrencia y modo de destino.}
\end{figure}
}{}

\IfFileExists{../stress/2026-07-05_local-smoke_v1/2026-07-05_fig-error-rate-vs-concurrency_v1.png}{
\begin{figure}[htbp]
  \centering
  \includegraphics[width=0.82\linewidth]{2026-07-05_fig-error-rate-vs-concurrency_v1.png}
  \caption{Tasa de error de la corrida local. No hubo conexiones fallidas.}
\end{figure}
}{}

\IfFileExists{../stress/2026-07-05_local-smoke_v1/2026-07-05_fig-latency-p50-p95-vs-concurrency_v1.png}{
\begin{figure}[htbp]
  \centering
  \includegraphics[width=0.82\linewidth]{2026-07-05_fig-latency-p50-p95-vs-concurrency_v1.png}
  \caption{Latencias p50/p95 totales de la corrida local.}
\end{figure}
}{}
